\section{Introduction}
\label{sec:intro}

Modern critical-infrastructure services run on open-source software.
Electronic-health-record platforms, financial-exchange back-ends,
telecommunications cores, government service portals, cloud platforms,
utility management systems, and logistics infrastructure are all
built, today, by assembling dozens of open-source components: database
engines, message brokers, container orchestrators, service meshes,
identity providers, caches, web servers, and custom glue code. Over
the last two decades, the dominant mode of building production
infrastructure has shifted from monolithic commercial software to
layered open-source stacks, and the services that depend on these
stacks have become increasingly networked, increasingly reachable
from the same internet that carries commodity cybercrime, and
increasingly fundamental to daily life. At the same
time, the consequences of a successful attack against them have grown
correspondingly severe: a compromise of a hospital records system,
a payment-clearing back-end, or a 911 call-taking platform is not
merely a data breach but a service-continuity incident with direct
harm to citizens. Recent years have made this concrete.
Ransomware has shut down hospital networks and payment rails for
weeks; attackers have pivoted through open-source build systems into
downstream production environments; and nation-state adversaries have
been observed pre-positioning inside U.S.\ critical-infrastructure
networks running commodity open-source stacks. The defensive side of
this arms race has not kept pace.

\subsection{The Problem: Why Assessing an Open-Source Stack Is Hard}
\label{sec:intro-problem}

Assessing the security posture of a specific deployment of an
open-source critical-infrastructure stack is categorically harder than
assessing a typical commercial-off-the-shelf application, for four
reinforcing reasons.

\paragraph{Heterogeneity and deployment-specific composition.}
A single \cps deployment typically contains a mixture of commodity
application servers, database engines (PostgreSQL, MySQL, MongoDB),
message brokers (Kafka, RabbitMQ, NATS), caches and key-value stores
(Redis, Memcached, etcd), container orchestrators (Kubernetes,
Nomad), service meshes and proxies (Envoy, HAProxy, nginx), identity
and policy services (Keycloak, OPA), coordination systems (Consul,
ZooKeeper), monitoring agents, and a substantial layer of custom,
site-specific code that stitches them together. Each of these
components has its own configuration surface, its own version
history, its own CVE backlog, and its own idiosyncratic interaction
with every other component it touches. Security tooling that
reasons about one component in isolation routinely misses
vulnerabilities that arise only from the specific combination of
versions and settings actually in use at a given site.

\paragraph{Operational sensitivity.}
The techniques that security teams routinely apply in isolated lab
environments---active network scanning, credential spraying, fuzzing,
exploit validation---cannot safely be used against a running \cps. An
aggressive scan can saturate a database connection pool and trigger
a cascading failure; a malformed request can crash a message broker
and lose in-flight events; a ``harmless'' probe of an authentication
endpoint can lock out legitimate users; fuzzing a production API
gateway can violate a regulated uptime SLA. This is not a hypothetical
concern: scan-induced outages and test-induced data corruption have
been documented repeatedly in production environments, and many
operators now contractually \emph{prohibit} active assessment of
production deployments. The result
is that the tests most likely to find real vulnerabilities are the ones
least likely to be permitted.

\paragraph{Dynamism.}
\CPSs are not static. Components are upgraded (or left un-upgraded),
services are added, removed, or relocated; third-party dependencies
come and go; feature flags change behavior seasonally or in response
to load; emergency patches and rollbacks leave behind
temporary-but-permanent configuration drift. An
assessment performed at time $t$ can be partially invalidated by time
$t + \Delta$, where $\Delta$ may be days, not years. Yet current
assessment cadences are typically annual, and sometimes much longer.

\paragraph{Scarce expertise.}
The specialists who can credibly evaluate such a system---engineers who
understand both the service semantics and the security-relevant software
stack---are few, expensive, and concentrated in a small number of
consultancies and national laboratories. No plausible growth in the
workforce will close this gap at the speed at which the attack surface
is growing.

\subsection{The State of the Art, and Why It Is Insufficient}
\label{sec:intro-sota}

Current practice for \cps security assessment reflects these
constraints. At the high end, a site will undergo a periodic manual
penetration test: a small team of experts spends days or weeks
interviewing operators, reading documentation, passively sniffing
traffic, and cautiously interacting with a limited subset of the system,
typically producing a report that is consumed once and then shelved. At
the lower end, passive asset inventory tools and network intrusion-detection
products observe network traffic and raise alerts against known-bad
patterns, but they do not reason about how a sophisticated adversary
might combine individually-innocuous weaknesses into an end-to-end
breach. Dependency-scanning and SBOM-driven tools flag known-vulnerable
component versions but do not reason about whether a particular CVE
is actually reachable in a given deployment, let alone how it composes
with other weaknesses. In between lie isolated-lab testbeds in which
an open-source component is analyzed in depth by its upstream community,
but such analyses inevitably diverge from any particular production
deployment, whose configuration, version mix, and custom glue code
are unique.

Three limitations are common to all of these approaches. First, they
are \emph{sporadic}: a site may be analyzed once a year, or once a
decade, while the threat landscape and the site itself evolve
continuously. Second, they are \emph{shallow with respect to
composition}: even excellent individual-component analyses stop short
of the multi-step attacks that characterize real incidents, because
such reasoning is performed (if at all) in the head of a human expert
whose time is limited. Third, they fail to exploit the central asset
of the defender---knowledge of the actual system's configuration,
documentation, and behavior---at scale, because there is no
infrastructure to translate that knowledge into automated adversarial
testing.

Meanwhile, attackers are moving in the opposite direction. Commodity
offensive tooling now includes AI-assisted reconnaissance, vulnerability
triage, and exploitation support. Recent evaluations of frontier
language models demonstrate measurable, rapid progress on multi-step
cyber-attack scenarios that only a few years ago would have required
a skilled human operator. The asymmetry is stark: the defender's
workflow is still artisanal, while the attacker's is increasingly
automated.

\subsection{The Proposed Approach: \sysname}
\label{sec:intro-approach}

We propose \sysname, an agentic framework for the continuous,
non-disruptive, in-depth security assessment of open-source
critical-infrastructure stacks.
The central idea is to move the adversarial activity off of the real
system and onto a \emph{faithful, continuously-updated digital twin},
where arbitrary---and otherwise unsafe---analyses can be performed
24~hours a day, 7~days a week, without any risk to the services the
real deployment provides.
Human experts remain in the loop, but their scarce attention is applied
where it matters most: confirming high-confidence findings, resolving
ambiguity that agents cannot, and validating suspected vulnerabilities
on the real system once a twin-based analysis has made the case that
validation is worth attempting.

\sysname is organized around four cooperating teams of AI agents, which
together form a closed feedback loop rather than a linear pipeline:

\begin{itemize}[leftmargin=1.3em,itemsep=0.15em,topsep=0.3em]
  \item \textbf{Knowledge Base (KB) Agents} ingest all forms of
  documentation about the target \cps---configuration files, logs,
  service manifests, network and deployment diagrams, software bills
  of materials, infrastructure-as-code artifacts, security policies,
  and change-management records---and organize this evidence into a hybrid
  knowledge base combining a graph database (for structured facts such
  as topology and component roles) and a vector database (for
  natural-language content that supports retrieval-augmented agent
  reasoning). KB Agents proactively interact with human operators when
  they detect missing or contradictory information.

  \item \textbf{Observer Agents} use the KB to construct an
  \emph{observation plan} that establishes strictly passive vantage
  points across the \cps: network taps, log readers, container-based
  sensors in specific subnets, service-mesh integrations, and eBPF
  probes where permitted. Observer Agents confirm or refute the KB's
  picture of the system and discover components, software versions,
  and flows that documentation missed. Their observations are
  continuous, so the KB evolves with the real system.

  \item \textbf{Replication Agents} consume the KB and observations to
  synthesize a digital twin: a virtualized, containerized, or
  process-level replica that reproduces the system's components, their
  software configuration, their network topology, and representative
  data. For components whose real deployment cannot be copied wholesale
  into the twin---proprietary sidecars, closed-source third-party
  services, licensed appliances interacting with the open-source stack,
  or components whose state volume makes full replication infeasible---
  Replication Agents build on the PIs' prior work on modeling opaque
  systems from observed interactions (Conware, HALucinator) and
  synthesize \emph{mock} components whose behavior is derived from
  observed API and network interactions. Crucially,
  fidelity is an explicit, quantified parameter of the replication
  algorithm rather than an implicit assumption.

  \item \textbf{Red Team Agents} continuously attack the digital twin.
  They compose \emph{component-level} analyses (fuzzing, static
  analysis, symbolic execution, taint tracking) with \emph{system-level}
  multi-step attack planning, building on our ChainReactor work on
  AI-planning-based privilege escalation and on the Artiphishell system
  developed for the DARPA AIxCC competition. Each candidate attack is
  grounded in a precondition/postcondition model expressed in PDDL, and
  its goal is to reach a nominated ``flag'' whose compromise has a
  clear security interpretation.
\end{itemize}

A finding produced by a Red Team Agent is not the end of the process.
Every confirmed attack on the twin is presented to a human security
analyst for triage. If the analyst judges the finding to be a
true-positive candidate, they validate it on the real system using an
appropriately cautious procedure, producing either a confirmed
vulnerability or additional evidence that the twin diverges from
reality. If the finding is judged to be a false positive caused by an
inaccurate twin, the fact of that inaccuracy is fed back to the KB,
Observer, and Replication Agents, which jointly plan updates to
documentation, observations, and configuration. Over time, this loop
drives both the twin's fidelity and the Red Team's precision upward.

\subsection{Contributions and Research Questions}
\label{sec:intro-contribs}

This project will produce the following concrete contributions.

\begin{enumerate}[leftmargin=1.5em,itemsep=0.15em,topsep=0.3em]
  \item A design for a \emph{hybrid, agent-native knowledge base} for
  open-source critical-infrastructure stacks that combines structured
  graph representations of topology and component roles with vectorized
  natural-language representations of documentation, policies, and
  observations, and that is resilient to partial, stale, and
  contradictory inputs.

  \item A suite of \emph{strictly-passive observation techniques} for
  \CPSs---network, log, API, configuration, and security-tool-based---
  and a principled method for synthesizing them into a continuously
  evolving system model.

  \item A \emph{controlled-fidelity digital twin synthesis algorithm}
  in which diversity of component classes, topological patterns, and
  data flows are explicit parameters, enabling minimally representative
  twins of large-scale systems.

  \item An \emph{agentic red-team workflow} that composes component-level
  vulnerability analyses into multi-step attack plans via PDDL-based
  planning, adapted from our ChainReactor and Artiphishell work and
  extended with the emergent reasoning capabilities of modern LLMs.

  \item An \emph{open-source reference implementation} of \sysname,
  together with a set of reference \cps testbeds and benchmarks that
  the research community can use to compare future approaches.
\end{enumerate}

These contributions will answer the following research questions:
\textbf{(RQ1)} Can a team of AI agents construct and maintain a useful,
living model of an open-source critical-infrastructure deployment from
the evidence that a typical operator already has?
\textbf{(RQ2)} What is the minimum fidelity a digital twin needs in
order to surface the same classes of vulnerabilities as the real
system, and how can that fidelity be measured?
\textbf{(RQ3)} Can AI-agent red teams, operating against a twin,
discover multi-step vulnerabilities with precision and recall
competitive with---and eventually exceeding---human red teams?
\textbf{(RQ4)} How should human experts be positioned in the loop so
that their scarce attention produces maximum marginal value?

\subsection{Organization of the Proposal}
The remainder of this proposal is organized as follows.
Section~\ref{sec:approach} describes the overall technical approach and
the four agent teams.
Section~\ref{sec:thrusts} details the three research thrusts that
constitute the bulk of the proposed work.
Section~\ref{sec:evaluation} presents our evaluation plan;
Section~\ref{sec:timeline} the project timeline;
Section~\ref{sec:broader} our broader impacts;
Section~\ref{sec:pis} the qualifications of the PIs;
and Section~\ref{sec:prior} results from prior NSF support.
